\chapter{Mạng nơ-ron thông tin vật lý}
\label{chap:pinn}
 
Ba chương trước đã dựng xong toàn bộ vật liệu. Chương~\ref{chap:fnn} cho một họ
hàm tham số hoá $\Net_{\btheta}$ trơn vô hạn theo biến đầu vào, cùng công thức
tường minh cho các đạo hàm ấy. Chương~\ref{chap:autodiff} cho cơ chế tính các đạo
hàm đó chính xác tới độ chính xác máy, cùng các thuật toán cực tiểu hoá. Chương
này ghép chúng lại thành một phương pháp giải phương trình vi phân, và -- quan
trọng không kém -- phân tích những chỗ phương pháp ấy hỏng.
 
Cấu trúc của chương theo ba lớp lồng nhau. Mục~\ref{sec:cau-truc-loss} xây dựng
phương pháp: phần dư, tập điểm phối trí, hàm mục tiêu. Trong đó,
Mục~\ref{subsec:chan-on-dinh} trả lời một câu hỏi mà tài liệu thường bỏ qua:
\emph{vì sao cực tiểu hoá phần dư lại là việc đáng làm?} Câu trả lời không hiển
nhiên -- phần dư nhỏ chỉ kéo theo sai số nhỏ khi bài toán vi phân có một chặn ổn
định, và hằng số của chặn ấy chính là đại lượng suy biến trong các bài toán khó.
Mục~\ref{sec:gradient-pathology} phân tích khó khăn ở tầng \emph{tối ưu}: nhiều
mục tiêu cạnh tranh trong một hàm mục tiêu duy nhất.
Mục~\ref{sec:spectral-bias} truy tiếp xuống tầng \emph{xấp xỉ}: ngay cả khi thuật
toán tối ưu hoàn hảo, mạng vẫn có thiên hướng nội tại không học được tần số cao.
Ba tầng này hội tụ về một kết luận định lượng thống nhất ở
Mục~\ref{subsec:he-qua-pinn}.
 
\subsubsection*{Quy ước ký hiệu bổ sung}
 
Chương này giữ Bảng~\ref{tab:ky-hieu} và bổ sung Bảng~\ref{tab:ky-hieu-4}. Hai
điểm cần lưu ý vì tài liệu gốc thường dùng khác.
 
Thứ nhất, $\Lop$ \emph{chỉ} ký hiệu toán tử vi phân, nhất quán từ
\eqref{eq:pinn-loss} của Chương~\ref{chap:fnn}; hàm mục tiêu và các thành phần
của nó dùng chữ $J$. Nhiều tài liệu, kể cả~\cite{wang2021}, dùng $\mathcal{L}$
cho hàm mất mát; ta tránh cách đó để không trùng ký hiệu.
 
Thứ hai, vectơ tham số vật lý được ký hiệu $\bzeta$ chứ không phải $\lambda$ như
trong~\cite{raissi2019}, bởi $\lambda$ trong báo cáo này được dành cho các trọng
số cân bằng mất mát. Hai đại lượng có bản chất hoàn toàn khác nhau -- một là dữ
kiện (hoặc ẩn số) của bài toán vật lý, một là siêu tham số của thuật toán -- nên
việc dùng chung ký hiệu là nguồn nhầm lẫn thực sự.
 
Thứ ba, tốc độ học được ký hiệu $\eta$ từ chương này trở đi, tương ứng với
$\alpha$ của \eqref{eq:iter-scheme} ở Chương~\ref{chap:autodiff}. Lý do đổi ký
hiệu: trong chương này $\alpha$ đã được dùng cho hệ số làm trơn của quy tắc cân
bằng trọng số \eqref{eq:lambda-update}, theo đúng ký hiệu của~\cite{wang2021}.

Để đọc thuận, ta viết $u_{\btheta} := \Net_{\btheta}$ khi mạng đóng vai trò hàm
nghiệm thử, và dành chữ $u$ không chỉ số cho nghiệm đúng.
 
\begin{table}[htbp]
\centering
\caption{Ký hiệu bổ sung của Chương~\ref{chap:pinn}.}
\label{tab:ky-hieu-4}
\small
\begin{tabular}{@{}l p{9.2cm}@{}}
\toprule
\textbf{Ký hiệu} & \textbf{Ý nghĩa} \\
\midrule
$\Omega \subset \R^{d}$, $\partial\Omega$ & miền không gian bị chặn và biên của
  nó \\
$\Lop[\,\cdot\,;\bzeta]$ & toán tử vi phân, phụ thuộc tham số vật lý $\bzeta$ \\
$\Bop$ & toán tử biên (Dirichlet, Neumann hoặc Robin) \\
$u$, $u_{\btheta} := \Net_{\btheta}$ & nghiệm đúng và nghiệm xấp xỉ bằng mạng \\
$r_{\btheta}$ & phần dư của phương trình, \eqref{eq:residual} \\
$e_{\btheta} := u_{\btheta} - u$ & sai số của nghiệm xấp xỉ \\
$\Tset_r, \Tset_0, \Tset_b$ & tập điểm phối trí trong miền, tại $t=0$, trên biên \\
$N_r, N_0, N_b$ & lực lượng của ba tập trên \\
$J_r, J_0, J_b$ & ba thành phần mất mát, \eqref{eq:loss-r}--\eqref{eq:loss-b} \\
$\lambda_r,\lambda_0,\lambda_b$ & trọng số cân bằng (siêu tham số của thuật
  toán) \\
$m$ & bậc cao nhất của đạo hàm xuất hiện trong $\Lop$ \\
\bottomrule
\end{tabular}
\end{table}
 
%==============================================================================
\section{Cấu trúc hàm mất mát}
\label{sec:cau-truc-loss}
%==============================================================================
 
\subsection{Bài toán tổng quát}
\label{subsec:bai-toan-tong-quat}
 
Xét bài toán giá trị đầu -- giá trị biên cho một phương trình đạo hàm riêng phi
tuyến có tham số, theo cách đặt vấn đề của Raissi và cộng sự~\cite{raissi2019}.
Với $\Omega \subset \R^{d}$ bị chặn và $[0,T]$ là khoảng thời gian khảo sát, ta
tìm $u : \Omega \times [0,T] \to \R$ thoả mãn
\begin{equation}
\frac{\partial u}{\partial t}(x,t) + \Lop\!\left[u;\, \bzeta\right](x,t) = 0,
\qquad (x,t) \in \Omega \times (0,T],
\label{eq:pde-general}
\end{equation}
kèm theo điều kiện đầu và điều kiện biên
\begin{align}
    u(x,0) &= u_{0}(x), && x \in \Omega, \label{eq:ic}\\
    \Bop\!\left[u\right](x,t) &= g(x,t),
      && (x,t) \in \partial\Omega \times (0,T].
    \label{eq:bc}
\end{align}
 
Ở đây $\Lop[\,\cdot\,;\bzeta]$ là toán tử vi phân theo biến không gian, có thể
phi tuyến, phụ thuộc vectơ tham số vật lý $\bzeta$; còn $\Bop$ là toán tử biên,
là toán tử đồng nhất (điều kiện Dirichlet), toán tử đạo hàm pháp tuyến (Neumann)
hoặc tổ hợp của chúng (Robin).
 
Dạng \eqref{eq:pde-general} bao hàm hầu hết các lớp phương trình được khảo sát
trong báo cáo. Với phương trình Burgers một chiều có độ nhớt $\nu$,
$\Lop[u] = u\,\partial_{x} u - \nu\, \partial_{xx} u$; với phương trình khuếch
tán, $\Lop[u] = -\nu\,\partial_{xx} u$. Bài toán dừng (không phụ thuộc thời gian)
thu được bằng cách bỏ số hạng $\partial_t u$ và điều kiện \eqref{eq:ic}; ta sẽ
dùng trường hợp này ở Mục~\ref{subsec:chan-on-dinh} vì nó cho phép phát biểu kết
quả gọn nhất.
 
Ta giả thiết xuyên suốt rằng bài toán \eqref{eq:pde-general}--\eqref{eq:bc}
\emph{đặt chỉnh}: nghiệm tồn tại, duy nhất và phụ thuộc liên tục vào dữ kiện. Giả
thiết này không phải hình thức. Mục~\ref{subsec:chan-on-dinh} sẽ cho thấy chính
mô-đun liên tục ấy -- chứ không phải bản thân sự tồn tại nghiệm -- là đại lượng
quyết định phương pháp có hoạt động hay không.
 
\subsection{Phần dư của phương trình}
\label{subsec:phan-du}
 
Ta xấp xỉ nghiệm $u$ bằng một mạng nơ-ron truyền thẳng theo
Định nghĩa~\ref{def:fnn}, với $n_0 = d+1$ (biến vào là $(x,t)$) và $n_L = 1$:
\begin{equation}
u_{\btheta} : \Omega \times [0,T] \to \R,
\qquad
(x,t) \longmapsto u_{\btheta}(x,t),
\qquad \btheta \in \R^{P}.
\label{eq:network-surrogate}
\end{equation}
 
Điểm mấu chốt phân biệt PINNs với các phương pháp học máy thuần dữ liệu là định
nghĩa sau.
 
\begin{definition}[Phần dư]
\label{def:residual}
Giả sử $u_{\btheta} \in C^{m}$ theo biến không gian và $C^1$ theo biến thời gian,
với $m$ là bậc cao nhất của đạo hàm trong $\Lop$. \emph{Phần dư} của phương
trình \eqref{eq:pde-general} tại điểm $(x,t)$ là
\begin{equation}
r_{\btheta}(x,t)
:= \frac{\partial u_{\btheta}}{\partial t}(x,t)
 + \Lop\!\left[u_{\btheta};\, \bzeta\right](x,t).
\label{eq:residual}
\end{equation}
\end{definition}
 
Ba tính chất của $r_{\btheta}$ cần được nhấn mạnh, mỗi tính chất nối với một kết
quả đã có ở các chương trước.
 
\emph{(a) Dùng chung tham số.} Hàm $r_{\btheta}$ dùng chung $\btheta$ với
$u_{\btheta}$: đây không phải một mạng nơ-ron thứ hai mà là kết quả của việc áp
toán tử vi phân lên chính đồ thị tính toán của $u_{\btheta}$
(Định nghĩa~\ref{def:comp-graph}). Do đó số tham số cần tối ưu \emph{không tăng}
khi ta đưa thông tin vật lý vào mô hình -- điều làm cho phương pháp khả thi, vì
theo Mệnh đề~\ref{md:chi-phi-bp}, chi phí gradient tỉ lệ với chi phí đánh giá hàm
chứ không với $P$.
 
\emph{(b) Đạo hàm theo biến đầu vào, tính bằng vi phân tự động.} Các đạo hàm
trong \eqref{eq:residual} là đạo hàm theo $(x,t)$, không phải theo $\btheta$.
Mệnh đề~\ref{md:ad-tai-tao} bảo đảm rằng vi phân tự động tính đúng các đại lượng
này, tái tạo chính xác hai hệ thức truy hồi \eqref{eq:G-truy-hoi} và
\eqref{eq:S-truy-hoi}. Khác với sai phân hữu hạn, cơ chế này không đòi hỏi lưới
rời rạc nào và không chịu rào cản độ chính xác của
Mệnh đề~\ref{md:fd-rao-can} -- một điểm mà Nhận xét~\ref{nx:fd-pde} đã lập luận
là điều kiện của phương pháp chứ không phải lựa chọn cài đặt.
 
\emph{(c) Ràng buộc lên hàm kích hoạt.} Giả thiết $u_{\btheta} \in C^m$ trong
Định nghĩa~\ref{def:residual} là ràng buộc trực tiếp lên kiến trúc, không phải
điều kiện kỹ thuật cho tiện. Theo Nhận xét~\ref{nx:bac-cao}, điều kiện cần để đạo
hàm bậc $m$ của mạng không đồng nhất triệt tiêu là $\sigma^{(m)} \not\equiv 0$;
và với ReLU, Hệ quả~\ref{hq:pinn-hong} cho kết luận mạnh hơn nhiều một sự ``suy
biến'': hàm mục tiêu trở thành \emph{hằng số địa phương}, nên mọi thuật toán
gradient đứng yên tuyệt đối. Hàm $\tanh$ thuộc $C^{\infty}(\R)$ với mọi đạo hàm
bị chặn đều (Hệ quả~\ref{hq:tanh-chan}), nên là lựa chọn đã được biện minh chứ
không phải quy ước; đó cũng là hàm kích hoạt dùng trong toàn bộ thực nghiệm ở
Chương~\ref{chap:thucnghiem}.
 
Toàn bộ luồng tính toán được tóm tắt trong Hình~\ref{fig:pinn-arch}.
 
\begin{figure}[htbp]
\centering
\resizebox{\textwidth}{!}{%
\begin{tikzpicture}[
  font=\small,
  neuron/.style={circle, draw=black!55, line width=0.5pt, minimum size=5.5mm, inner sep=0pt, fill=white},
  box/.style={rounded corners=2pt, draw=black!60, line width=0.6pt, align=center, inner sep=4.5pt, fill=white},
  lossbox/.style={box, fill=blue!6},
  flow/.style={-{Stealth[length=4.5pt]}, line width=0.7pt, black!70},
  wire/.style={draw=black!20, line width=0.28pt},
  grad/.style={-{Stealth[length=5pt]}, line width=0.8pt, dashed, red!62!black}
]

\node[box, fill=orange!10] (in) at (0,0) {$(x,t)$};
\def\lxA{2.0} \def\lxB{3.3} \def\lxC{4.6} \def\lxD{5.9}

\foreach \i in {1,2} \node[neuron] (a\i) at (\lxA, {0.75-1.5*(\i-1)}) {};
\foreach \i in {1,...,4} \node[neuron] (b\i) at (\lxB, {1.35-0.9*(\i-1)}) {};
\foreach \i in {1,...,4} \node[neuron] (c\i) at (\lxC, {1.35-0.9*(\i-1)}) {};
\node[neuron] (d1) at (\lxD, 0) {};

\begin{scope}[on background layer]
  \foreach \i in {1,2} \foreach \j in {1,...,4} \draw[wire] (a\i)--(b\j);
  \foreach \i in {1,...,4} \foreach \j in {1,...,4} \draw[wire] (b\i)--(c\j);
  \foreach \i in {1,...,4} \draw[wire] (c\i)--(d1);
\end{scope}

\node[fit=(a1)(a2)(b1)(b4)(c1)(c4)(d1), inner sep=3pt, draw=none] (netbox) {};
\node[font=\footnotesize, black!65] at ($(netbox.north)+(1.55,0.30)$) {mạng truyền thẳng $u_{\btheta}$};

\draw[flow] (in) -- (a1);
\draw[flow] (in) -- (a2);

\node[box, fill=orange!10] (out) at (7.5,0) {$u_{\btheta}(x,t)$};
\draw[flow] (d1) -- (out);

\node[box, fill=green!7] (ad) at (10.6,0) {vi phân tự động\\[2pt] $\dfrac{\partial u_{\btheta}}{\partial t},\; \dfrac{\partial u_{\btheta}}{\partial x},\; \dfrac{\partial^{2} u_{\btheta}}{\partial x^{2}}$};
\draw[flow] (out) -- (ad);

\node[box, fill=green!7] (res) at (14.4,0) {phần dư\\[3pt] $r_{\btheta} = \dfrac{\partial u_{\btheta}}{\partial t} + \Lop[u_{\btheta};\bzeta]$};
\draw[flow] (ad) -- (res);

\node[lossbox] (lb) at (1.7,-3.5) {$J_b=\dfrac{1}{N_b}\displaystyle\sum |\Bop[u_{\btheta}]-g|^{2}$};
\node[lossbox] (l0) at (7.5,-3.5) {$J_0=\dfrac{1}{N_0}\displaystyle\sum |u_{\btheta}-u_0|^{2}$};
\node[lossbox] (lr) at (13.2,-3.5) {$J_r=\dfrac{1}{N_r}\displaystyle\sum |r_{\btheta}|^{2}$};

\draw[flow] (res.south) -- ++(0,-0.55) -| (lr.north);
\draw[flow] (out.south) -- (l0.north);
\draw[flow] (out.south) -- ++(0,-0.9) -| (lb.north);

\node[box, fill=blue!14, line width=0.85pt] (ltot) at (7.5,-5.6) {$J(\btheta)=\lambda_r J_r+\lambda_0 J_0+\lambda_b J_b$};

\draw[flow] (lb.south) |- ($(ltot.west)+(-0.9,0)$) -- (ltot.west);
\draw[flow] (lr.south) |- ($(ltot.east)+(0.9,0)$) -- (ltot.east);
\draw[flow] (l0.south) -- (ltot.north);

\draw[grad] (ltot.west) -- ++(-1.6,0) -- (-1.5,-5.6) -- (-1.5,2.45) -- (2.0,2.45) -- (2.0,1.18);

\node[font=\footnotesize, red!62!black, align=center] at (-0.55,-4.35) {cập nhật tham số\\[-1pt] $\nabla_{\btheta} J$};
\end{tikzpicture}}
\caption[Luồng tính toán của một bước huấn luyện PINN]{Luồng tính toán của một
bước huấn luyện PINN. Mạng truyền thẳng $u_{\btheta}$ nhận $(x,t)$ và trả về giá
trị nghiệm thử; vi phân tự động (Chương~\ref{chap:autodiff}) lấy các đạo hàm
theo \emph{biến đầu vào} để dựng phần dư \eqref{eq:residual}. Ba thành phần mất
mát \eqref{eq:loss-r}--\eqref{eq:loss-b} được gộp theo trọng số thành
\eqref{eq:total-loss}; mũi tên đứt màu đỏ là lượt ngược lấy
$\nabla_{\btheta}J$ trên \emph{toàn bộ} đồ thị mở rộng, tức đi xuyên qua cả khối
vi phân tự động. Điểm cần chú ý: chỉ có \emph{một} bộ tham số $\btheta$ duy nhất
cho cả ba nhánh mất mát.}
\label{fig:pinn-arch}
\end{figure}
 
\subsection{Điểm phối trí và chiến lược lấy mẫu}
\label{subsec:collocation}
 
Để đánh giá phần dư một cách số trị, ta chọn một tập hữu hạn điểm trong miền tính
toán. Theo thuật ngữ của phương pháp trọng số dư, các điểm này được gọi là
\emph{điểm phối trí} (collocation point). Ta xây dựng ba tập:
\begin{align}
\Tset_{r} &= \{(x_{r}^{i}, t_{r}^{i})\}_{i=1}^{N_{r}}
  \subset \Omega \times (0,T], && \text{điểm trong miền,} \label{eq:Tr}\\
\Tset_{0} &= \{x_{0}^{i}\}_{i=1}^{N_{0}} \subset \Omega,
  && \text{điểm trên lát cắt } t=0, \label{eq:T0}\\
\Tset_{b} &= \{(x_{b}^{i}, t_{b}^{i})\}_{i=1}^{N_{b}}
  \subset \partial\Omega \times (0,T], && \text{điểm trên biên.} \label{eq:Tb}
\end{align}
 
Đặc điểm đáng chú ý của $\Tset_{r}$ là các điểm này \emph{không cần nhãn}. Ta
không biết và không cần biết giá trị nghiệm đúng tại chúng; thông tin giám sát
duy nhất là điều kiện $r_{\btheta} = 0$ do chính phương trình áp đặt. Đây là
nguồn gốc của tên gọi ``thông tin vật lý'': phương trình đóng vai trò thay thế
cho dữ liệu nhãn. Ngược lại, $\Tset_0$ và $\Tset_b$ \emph{có} nhãn, lấy từ
$u_0$ và $g$.
 
Về cách lấy mẫu, Raissi và cộng sự~\cite{raissi2019} dùng \emph{lấy mẫu siêu khối
Latin} (Latin Hypercube Sampling, LHS). So với lấy mẫu ngẫu nhiên đều, LHS bảo
đảm mỗi chiều của miền được phủ đồng đều hơn khi $N_{r}$ còn khiêm tốn, nhờ đó
giảm phương sai của ước lượng phần dư ở \eqref{eq:mc-estimate}. Các dãy tựa ngẫu
nhiên khác như Sobol hay Halton cho kết quả tương đương. Cần nói rõ rằng lợi thế
này là lợi thế \emph{phương sai hữu hạn mẫu}: cả ba chiến lược đều cho cùng một
giới hạn khi $N_r \to \infty$, và khác biệt mờ dần khi $N_r$ tăng. Ảnh hưởng của
chiến lược lấy mẫu đến độ chính xác sẽ được so sánh trong phần thực nghiệm.
 
\begin{figure}[htbp]
\centering
\includegraphics[width=\textwidth]{Images/chap_4/hinh_collocation.pdf}
\caption[Phân bố các tập điểm huấn luyện của PINNs]{Phân bố ba tập điểm huấn
luyện trên miền $[-1,1]\times[0,1]$. Bảng (a) dùng lấy mẫu siêu khối Latin, bảng
(b) dùng lấy mẫu ngẫu nhiên đều, cùng số điểm $N_r = 150$. Lưới nền chia miền
thành $10\times 10$ ô nhằm làm rõ mức độ phủ: với LHS, hình chiếu của tập điểm
lên mỗi trục toạ độ phủ đều các khoảng chia, trong khi lấy mẫu đều để lại những ô
trống và những cụm điểm dày. Khác biệt này rõ nhất khi $N_r$ còn nhỏ và mờ dần
khi $N_r$ tăng.}
\label{fig:collocation}
\end{figure}
 
\subsection{Hàm mất mát tổng hợp}
\label{subsec:loss}
 
Với ba tập điểm nói trên, ta định nghĩa ba thành phần mất mát dưới dạng sai số
bình phương trung bình:
\begin{align}
J_{r}(\btheta)
  &= \frac{1}{N_{r}} \sum_{i=1}^{N_{r}}
     \abs{r_{\btheta}\!\left(x_{r}^{i}, t_{r}^{i}\right)}^{2},
     \label{eq:loss-r}\\[2pt]
J_{0}(\btheta)
  &= \frac{1}{N_{0}} \sum_{i=1}^{N_{0}}
     \abs{u_{\btheta}\!\left(x_{0}^{i}, 0\right)
          - u_{0}\!\left(x_{0}^{i}\right)}^{2},
     \label{eq:loss-0}\\[2pt]
J_{b}(\btheta)
  &= \frac{1}{N_{b}} \sum_{i=1}^{N_{b}}
     \abs{\Bop\!\left[u_{\btheta}\right]\!\left(x_{b}^{i}, t_{b}^{i}\right)
          - g\!\left(x_{b}^{i}, t_{b}^{i}\right)}^{2}.
     \label{eq:loss-b}
\end{align}
Việc chọn dạng bình phương không tuỳ tiện: theo
Mệnh đề~\ref{md:gauss-mse}, sai số bình phương trung bình là hệ quả của nguyên lý
hợp lý cực đại dưới giả thiết nhiễu Gauss đồng phương sai. Ở đây giả thiết ấy
không có nội dung xác suất thực sự cho $J_r$ -- phần dư không phải một quan trắc
có nhiễu -- nhưng dạng bình phương vẫn được biện minh bởi
Nhận xét~\ref{rem:monte-carlo} dưới đây theo một cách khác hẳn.
 
Hàm mất mát tổng hợp là tổ hợp tuyến tính có trọng số:
\begin{equation}
\boxed{\;
J(\btheta)
= \lambda_{r}\, J_{r}(\btheta)
+ \lambda_{0}\, J_{0}(\btheta)
+ \lambda_{b}\, J_{b}(\btheta)
\;}
\label{eq:total-loss}
\end{equation}
với $\lambda_{r}, \lambda_{0}, \lambda_{b} > 0$. Bài toán huấn luyện trở thành bài
toán tối ưu không ràng buộc
\begin{equation}
\btheta^{\star} \in \arg\min_{\btheta \in \R^{P}} J(\btheta),
\label{eq:opt-problem}
\end{equation}
giải bằng các thuật toán ở Mục~\ref{sec:optimizers}. Ta viết $\in$ chứ không viết
$=$ vì theo Mệnh đề~\ref{md:khong-loi} và Nhận xét~\ref{nx:doi-xung-hoan-vi}, tập
cực tiểu không phải một điểm mà chứa ít nhất $2^{n_1}n_1!$ phần tử tương đương.
 
\begin{nhanxet}[Bản chất ước lượng Monte Carlo của $J_{r}$]
\label{rem:monte-carlo}
Thành phần $J_{r}$ trong \eqref{eq:loss-r} không phải một hàm mất mát giám sát
thông thường. Nếu các điểm phối trí được lấy mẫu độc lập theo phân phối đều trên
$\Omega \times (0,T]$ và $r_{\btheta} \in L^2$, thì theo luật số lớn mạnh
\begin{equation}
J_{r}(\btheta)
\;\xrightarrow[\;N_{r} \to \infty\;]{\text{h.c.c.}}\;
\frac{1}{\abs{\Omega}\,T}
\int_{0}^{T}\!\!\int_{\Omega}
\abs{r_{\btheta}(x,t)}^{2}\, dx\, dt
= \frac{1}{\abs{\Omega}\,T}\,
\norm{r_{\btheta}}_{L^{2}(\Omega \times (0,T))}^{2},
\label{eq:mc-estimate}
\end{equation}
và với mỗi $\btheta$ cố định, $J_r(\btheta)$ là ước lượng \emph{không chệch} của
vế phải. Cách nhìn này cho thấy PINNs thực chất là một phương pháp \emph{bình
phương tối thiểu trong không gian hàm}, với không gian thử là tập các hàm biểu
diễn được bởi kiến trúc đã chọn,
$\mathcal{U}_{\btheta} = \{u_{\btheta} : \btheta \in \R^{P}\}$. Sai số tổng thể
do đó tách thành ba nguồn: sai số xấp xỉ (do $\mathcal{U}_{\btheta}$ không chứa
nghiệm đúng -- Định lý~\ref{dl:sobolev} bảo đảm nguồn này nhỏ tuỳ ý), sai số lấy
mẫu (do $N_{r}$ hữu hạn) và sai số tối ưu (do thuật toán không đạt cực tiểu toàn
cục).
 
Một dè dặt quan trọng thường bị bỏ qua. Tính không chệch ở trên đòi hỏi $\btheta$
\emph{độc lập} với mẫu $\Tset_r$. Trong thực hành, $\Tset_r$ được lấy một lần rồi
cố định, và $\btheta^{\star}$ được tối ưu \emph{trên chính mẫu đó}; khi ấy
$J_r(\btheta^{\star})$ là ước lượng chệch xuống của
$\norm{r_{\btheta^\star}}^2_{L^2}$ -- đúng hiện tượng quá khớp quen thuộc, ở đây
là quá khớp vào tập điểm phối trí. Hệ quả thực hành: giá trị $J_r$ cuối cùng
\emph{không} phải một ước lượng đáng tin cho chuẩn phần dư; phải đánh giá lại
$r_{\btheta^\star}$ trên một tập điểm độc lập. Ta áp dụng nguyên tắc này trong
toàn bộ thực nghiệm ở Chương~\ref{chap:thucnghiem}. Cần lưu ý rằng chính việc cố
định $\Tset_r$ lại là điều kiện cho phép dùng L-BFGS
(Mục~\ref{subsec:adam-vs-lbfgs}): hai yêu cầu này xung đột nhau, và ta chọn giữ
tập cố định rồi kiểm tra lại trên tập độc lập.
\end{nhanxet}
 
\subsection{Vì sao cực tiểu hoá phần dư là việc đáng làm}
\label{subsec:chan-on-dinh}
 
Toàn bộ phương pháp đứng trên một suy luận chưa được kiểm chứng: ta cực tiểu hoá
$\norm{r_{\btheta}}$ nhưng cái ta thực sự quan tâm là $\norm{e_{\btheta}}$ với
$e_{\btheta} = u_{\btheta} - u$. Hai đại lượng này không đồng nhất, và không có
lý do tiên nghiệm nào bảo đảm phần dư nhỏ kéo theo sai số nhỏ. Mối liên hệ giữa
chúng là một \emph{chặn ổn định} của bài toán vi phân, và mục này thiết lập nó
cho trường hợp đơn giản nhất -- đủ để thấy hằng số của chặn phụ thuộc vào cái gì.
 
\begin{menhde}[Chặn sai số qua phần dư: bài toán Poisson một chiều]
\label{md:chan-on-dinh}
Xét $-u'' = f$ trên $(0,1)$ với $u(0) = u(1) = 0$, và $u_{\btheta} \in C^2$ thoả
mãn \emph{chính xác} điều kiện biên. Đặt $r_{\btheta} := -u_{\btheta}'' - f$ và
$e_{\btheta} := u_{\btheta} - u$. Khi đó
\begin{equation}
\norm{e_{\btheta}'}_{L^2(0,1)} \;\le\; \frac{1}{\pi}\,
  \norm{r_{\btheta}}_{L^2(0,1)},
\qquad
\norm{e_{\btheta}}_{L^2(0,1)} \;\le\; \frac{1}{\pi^{2}}\,
  \norm{r_{\btheta}}_{L^2(0,1)} .
\label{eq:chan-poisson}
\end{equation}
\end{menhde}
 
\begin{proof}
Trừ hai phương trình: $-e_{\btheta}'' = -u_{\btheta}'' + u'' = -u_{\btheta}'' - f
= r_{\btheta}$, với $e_{\btheta}(0) = e_{\btheta}(1) = 0$. Nhân hai vế với
$e_{\btheta}$ và lấy tích phân trên $(0,1)$; tích phân từng phần số hạng vế trái,
số hạng biên triệt tiêu do $e_{\btheta}$ triệt tiêu tại hai đầu mút:
\[
\int_0^1 \big(e_{\btheta}'\big)^2\,dx
 = \int_0^1 r_{\btheta}\, e_{\btheta}\,dx
 \;\le\; \norm{r_{\btheta}}_{L^2}\,\norm{e_{\btheta}}_{L^2},
\]
bất đẳng thức cuối là Cauchy--Schwarz. Bất đẳng thức Poincaré trên $(0,1)$ với
điều kiện biên thuần nhất cho $\norm{e_{\btheta}}_{L^2} \le
\tfrac{1}{\pi}\norm{e_{\btheta}'}_{L^2}$, trong đó hằng số $1/\pi$ là tối ưu (đạt
tại $e_{\btheta}(x) = \sin\pi x$, hàm riêng ứng với trị riêng nhỏ nhất $\pi^2$
của $-d^2/dx^2$). Thay vào:
\[
\norm{e_{\btheta}'}^2_{L^2}
 \le \norm{r_{\btheta}}_{L^2}\cdot\tfrac{1}{\pi}\norm{e_{\btheta}'}_{L^2},
\]
chia hai vế cho $\norm{e_{\btheta}'}_{L^2}$ (trường hợp bằng $0$ là tầm thường)
được bất đẳng thức thứ nhất; áp Poincaré lần nữa được bất đẳng thức thứ hai.
\end{proof}
 
Mệnh đề~\ref{md:chan-on-dinh} là lời biện minh mà phương pháp cần: nó biến việc
cực tiểu hoá một đại lượng \emph{tính được} ($r_{\btheta}$, không cần biết
nghiệm) thành một chặn cho đại lượng \emph{quan tâm} ($e_{\btheta}$, không tính
được). Hằng số $1/\pi^2 \approx 0{,}101$ thậm chí có lợi: phần dư bị \emph{giảm}
đi khi chuyển sang sai số.
 
Kết quả sau cho thấy tình hình đảo ngược hoàn toàn với bài toán nhiễu kỳ dị.
 
\begin{menhde}[Chặn ổn định suy biến với bài toán đối lưu--khuếch tán]
\label{md:chan-suy-bien}
Xét $-\nu u'' + u' = f$ trên $(0,1)$ với $u(0) = u(1) = 0$ và $\nu > 0$, và
$u_{\btheta} \in C^2$ thoả mãn chính xác điều kiện biên. Với
$r_{\btheta} := -\nu u_{\btheta}'' + u_{\btheta}' - f$,
\begin{equation}
\norm{e_{\btheta}'}_{L^2} \le \frac{1}{\pi\,\nu}\,\norm{r_{\btheta}}_{L^2},
\qquad
\norm{e_{\btheta}}_{L^2} \le \frac{1}{\pi^{2}\,\nu}\,\norm{r_{\btheta}}_{L^2}.
\label{eq:chan-suy-bien}
\end{equation}
\end{menhde}
 
\begin{proof}
Như trên, $-\nu e_{\btheta}'' + e_{\btheta}' = r_{\btheta}$ với
$e_{\btheta}(0) = e_{\btheta}(1) = 0$. Nhân với $e_{\btheta}$, lấy tích phân và
tích phân từng phần:
\[
\nu\int_0^1 \big(e_{\btheta}'\big)^2 dx
 + \int_0^1 e_{\btheta}'\,e_{\btheta}\,dx
 = \int_0^1 r_{\btheta}\,e_{\btheta}\,dx .
\]
Số hạng đối lưu triệt tiêu:
$\int_0^1 e_{\btheta}'e_{\btheta}\,dx
= \big[\tfrac12 e_{\btheta}^2\big]_0^1 = 0$ do điều kiện biên. Phần còn lại lặp
lại đúng lập luận Cauchy--Schwarz và Poincaré của
Mệnh đề~\ref{md:chan-on-dinh}, với thừa số $\nu$ ở vế trái.
\end{proof}
 
\begin{nhanxet}[Đọc hai mệnh đề trên cùng nhau]
\label{nx:doc-chan}
So sánh \eqref{eq:chan-poisson} với \eqref{eq:chan-suy-bien}: hằng số ổn định
tăng theo $1/\nu$. Với $\nu = 0{,}01/\pi \approx 3{,}18\times10^{-3}$ -- giá trị
độ nhớt dùng cho phương trình Burgers ở Mục~\ref{sec:spectral-bias} -- thừa số
khuếch đại là $1/(\pi^2\nu) \approx 31{,}8$. Nói cách khác, để đạt cùng một sai
số $\norm{e_{\btheta}}$, bài toán nhiễu kỳ dị đòi hỏi phần dư nhỏ hơn khoảng ba
trăm lần so với bài toán Poisson.
 
Đây là điểm cần nhấn mạnh vì nó thường bị lẫn với thiên kiến phổ. Ngay cả với một
thuật toán tối ưu hoàn hảo và một lớp hàm biểu diễn được nghiệm, việc đưa hàm mục
tiêu xuống một mức $J_r$ cho trước \emph{không} bảo đảm sai số nhỏ như nhau giữa
hai bài toán. Mục~\ref{subsec:he-qua-pinn} sẽ cho thấy tính chất này cộng hưởng
với thiên kiến phổ theo chiều bất lợi.
 
Ba giới hạn của hai mệnh đề trên cần nêu rõ. Thứ nhất, cả hai giả thiết điều kiện
biên được thoả \emph{chính xác}; với ràng buộc mềm
(Nhận xét~\ref{rem:soft-constraint}), vế phải phải cộng thêm một số hạng tỉ lệ
với sai số biên. Thứ hai, cả hai là bài toán tuyến tính một chiều; với toán tử
phi tuyến như Burgers, chặn tương ứng chỉ tồn tại địa phương và hằng số phụ thuộc
nghiệm. Thứ ba, chúng chặn sai số theo chuẩn $L^2$, còn cái ta thường báo cáo là
sai số tương đối theo chuẩn vô cùng. Không kết quả nào trong báo cáo cho phép
chuyển đổi tự động giữa các chuẩn ấy; ta dẫn hai mệnh đề như \emph{cơ sở khái
niệm} cho phương pháp và như một dự báo định lượng về vai trò của $\nu$, không
như một chặn hậu nghiệm áp dụng được cho các thí nghiệm.
\end{nhanxet}
 
\subsection{Ràng buộc mềm và ràng buộc cứng}
\label{subsec:rang-buoc}
 
\begin{nhanxet}[Ràng buộc mềm và hệ quả]
\label{rem:soft-constraint}
Công thức \eqref{eq:total-loss} áp đặt điều kiện đầu và điều kiện biên dưới dạng
\emph{ràng buộc mềm}: chúng chỉ được phạt trong hàm mục tiêu chứ không được thoả
mãn chính xác. Đây là khác biệt căn bản so với phương pháp phần tử hữu hạn, nơi
điều kiện Dirichlet được áp cứng lên không gian thử bằng cách chọn cơ sở triệt
tiêu trên biên.
 
Hai hệ quả. Thứ nhất, nghiệm số thu được nói chung \emph{không} thoả mãn chính
xác điều kiện biên, nên giả thiết của Mệnh đề~\ref{md:chan-on-dinh} không thoả
và chặn sai số phải được bổ sung số hạng biên. Thứ hai, và nghiêm trọng hơn, việc
gộp nhiều mục tiêu cạnh tranh vào một hàm mục tiêu duy nhất tạo ra chính vấn đề
được phân tích ở Mục~\ref{sec:gradient-pathology}.
\end{nhanxet}
 
Hướng khắc phục triệt để là dựng nghiệm thử thoả mãn điều kiện biên theo cấu
trúc, tức \emph{ràng buộc cứng}.
 
\begin{bode}[Nghiệm thử với ràng buộc cứng]
\label{bd:rang-buoc-cung}
Cho $G \in C^{m}$ thoả $G|_{\partial\Omega} = g$ và $D \in C^{m}$ thoả
$D|_{\partial\Omega} = 0$. Đặt
\begin{equation}
\hat{u}_{\btheta}(x,t) := G(x,t) + D(x,t)\, u_{\btheta}(x,t).
\label{eq:hard-constraint}
\end{equation}
Khi đó $\hat{u}_{\btheta}$ thoả mãn điều kiện biên Dirichlet \eqref{eq:bc}
\emph{với mọi} $\btheta \in \R^P$, nên $J_b \equiv 0$ và thành phần này bị loại
khỏi \eqref{eq:total-loss}.
\end{bode}
 
\begin{proof}
Với $(x,t) \in \partial\Omega\times(0,T]$ ta có $D(x,t) = 0$, nên
$\hat u_{\btheta}(x,t) = G(x,t) = g(x,t)$ bất kể giá trị $u_{\btheta}(x,t)$.
Tính trơn của $\hat u_{\btheta}$ suy từ tính trơn của $G$, $D$ và
$u_{\btheta}$ cùng quy tắc tích.
\end{proof}
 
Bổ đề~\ref{bd:rang-buoc-cung} loại bỏ hoàn toàn một mục tiêu cạnh tranh và khôi
phục giả thiết của Mệnh đề~\ref{md:chan-on-dinh}. Cái giá là việc dựng $G$ và $D$
tường minh chỉ khả thi với hình học đơn giản: trên $\Omega = (a,b)$ ta có thể lấy
$D(x) = (x-a)(b-x)$, nhưng trên một miền hai chiều có góc hoặc biên cong, việc
tìm $D$ trơn triệt tiêu đúng trên $\partial\Omega$ và khác không bên trong đã là
một bài toán riêng. Ngoài ra, \eqref{eq:hard-constraint} thay đổi lớp hàm thử: độ
trơn và thang độ của $\hat u_{\btheta}$ bị chi phối bởi $D$, nên các phân tích ở
Chương~\ref{chap:fnn} về $u_{\btheta}$ không chuyển sang $\hat u_{\btheta}$ một
cách tự động. Báo cáo dùng ràng buộc mềm cho các thí nghiệm chính và ràng buộc
cứng như một đối chứng.
 
\subsection{Bài toán thuận và bài toán ngược}
\label{subsec:thuan-nguoc}
 
Khung hàm mất mát \eqref{eq:total-loss} bao trùm cả hai lớp bài toán được Raissi
và cộng sự~\cite{raissi2019} khảo sát.
 
\paragraph{Bài toán thuận.}
Tham số vật lý $\bzeta$ được cho trước; ẩn số duy nhất là $\btheta$. Bài toán tối
ưu chính là \eqref{eq:opt-problem}. Toàn bộ thông tin giám sát nằm ở điều kiện
đầu, điều kiện biên và phương trình. Đây là trường hợp của TN1--TN10 ở
Chương~\ref{chap:thucnghiem}.
 
\paragraph{Bài toán ngược.}
Tham số $\bzeta$ chưa biết và cần được xác định từ một tập quan trắc rời rạc
$\{(x^{i}, t^{i}, u^{i})\}_{i=1}^{N_{u}}$ nằm rải rác trong miền, chứ không chỉ
trên biên. Hàm mục tiêu trở thành
\begin{equation}
J(\btheta, \bzeta)
= \underbrace{\frac{1}{N_{u}} \sum_{i=1}^{N_{u}}
   \abs{u_{\btheta}(x^{i}, t^{i}) - u^{i}}^{2}}_{\text{khớp dữ liệu quan trắc}}
+ \underbrace{\frac{1}{N_{r}} \sum_{j=1}^{N_{r}}
   \abs{r_{\btheta,\bzeta}(x_{r}^{j}, t_{r}^{j})}^{2}}_{\text{ràng buộc vật lý}},
\label{eq:inverse-loss}
\end{equation}
và ta tối ưu \emph{đồng thời} theo cặp $(\btheta, \bzeta)$. Điểm đáng chú ý là
$\bzeta$ được xử lý hoàn toàn giống một tham số của mạng: theo
Định nghĩa~\ref{def:comp-graph}, nó là một đỉnh nguồn trong đồ thị tính toán và
nhận gradient qua cùng cơ chế \eqref{eq:chain-graph}, không cần một thuật toán
riêng. Phương pháp cổ điển thì cần đúng một thuật toán riêng: rời rạc hoá bài
toán thuận, rồi suy ra và cài đặt bài toán liên hợp của nó để lấy gradient. Bản
trước của báo cáo gọi điểm khác biệt này là ``ưu thế nổi bật''.
Mục~\ref{sec:tn11} đặt hai cách làm cạnh nhau với cùng dữ liệu và thấy ưu thế ấy
có thật nhưng hẹp hơn: nó nằm ở \emph{công sức xây dựng} phương pháp, còn về độ
chính xác và thời gian chạy thì phương pháp liên hợp ngang ngửa, thậm chí nhanh
hơn (Nhận xét~\ref{nx:doc-tn11}).
 
\begin{vidu}[Phát hiện tham số của phương trình Burgers]
\label{vd:burgers-inverse}
Xét phương trình Burgers dưới dạng tham số hoá
\begin{equation}
\frac{\partial u}{\partial t}
+ \zeta_{1}\, u\, \frac{\partial u}{\partial x}
- \zeta_{2}\, \frac{\partial^{2} u}{\partial x^{2}} = 0,
\label{eq:burgers-param}
\end{equation}
với giá trị đúng $\zeta_{1} = 1$ và
$\zeta_{2} = 0{,}01/\pi \approx 3{,}1831\times 10^{-3}$. Raissi và cộng
sự~\cite{raissi2019} huấn luyện trên $N_u = 2000$ điểm quan trắc rải trong miền
không gian--thời gian, với mạng chín lớp --- tám lớp ẩn, $20$ nơ-ron mỗi lớp
ẩn, theo đúng mã nguồn công bố kèm công trình. Kết quả nhận dạng: với dữ liệu sạch, $\zeta_1 = 0{,}99915$ và
$\zeta_2 = 3{,}1794\times10^{-3}$, tức sai số tương đối $0{,}09\%$ và
$0{,}12\%$; với dữ liệu nhiễu Gauss $1\%$, $\zeta_1 = 1{,}00042$ và
$\zeta_2 = 3{,}2098\times10^{-3}$, tức $0{,}04\%$ và $0{,}84\%$.
 
Hai kết luận, và một dè dặt. Kết luận thứ nhất: ràng buộc vật lý đóng vai trò
chính quy hoá mạnh, giúp bài toán ngược vốn đặt không chỉnh trở nên ổn định.
Kết luận thứ hai: sai số của $\zeta_2$ \emph{thường} lớn hơn sai số của
$\zeta_1$, và điều này có lý do -- $\zeta_2$ nhân với $\partial_{xx}u$, đại
lượng khó ước lượng nhất từ dữ liệu, đúng theo tinh thần
Nhận xét~\ref{nx:fd-pde}. Bản trước của báo cáo viết ``luôn lớn hơn khoảng một
bậc''; Mục~\ref{sec:tn11} lặp thí nghiệm trên năm hạt giống và cho thấy điều ấy
đúng về trung vị (tỉ số $24{,}5$) nhưng chỉ ở tám trên mười lần chạy --- và
chính cặp số sạch ở trên, $0{,}09\%$ với $0{,}12\%$, vốn đã không ủng hộ chữ
``luôn''. Dè dặt: các bảng khảo sát độ nhạy
trong~\cite{raissi2019} cho thấy độ chính xác suy giảm nhanh khi mức nhiễu tăng
lên $5\%$ hoặc $10\%$, và khi mạng quá nông; con số dưới $1\%$ ở trên đặc trưng
cho cấu hình cụ thể vừa nêu, không phải cho phương pháp nói chung.
\end{vidu}
 
%==============================================================================
\section{Không gian hàm mất mát và bệnh lý gradient}
\label{sec:gradient-pathology}
%==============================================================================
 
Mục trước xây dựng hàm mất mát \eqref{eq:total-loss} như một tổng có trọng số.
Cách đặt vấn đề này đơn giản về hình thức nhưng ẩn chứa một khó khăn tối ưu
nghiêm trọng, được Wang, Teng và Perdikaris~\cite{wang2021} phân tích hệ thống
dưới tên gọi \emph{bệnh lý gradient} (gradient pathology).
 
\subsection{Sự cạnh tranh giữa các thành phần mất mát}
\label{subsec:canh-tranh}
 
Để thuận tiện, ta viết lại \eqref{eq:total-loss} dưới dạng thu gọn
\begin{equation}
J(\btheta)
= J_{r}(\btheta)
+ \sum_{i=1}^{M} \lambda_{i}\, J_{i}(\btheta),
\label{eq:loss-compact}
\end{equation}
trong đó ta đã chuẩn hoá $\lambda_{r} = 1$ -- hợp lệ vì \eqref{eq:opt-problem}
bất biến dưới phép nhân $J$ với một hằng số dương -- và gộp các thành phần điều
kiện đầu, điều kiện biên vào họ $\{J_{i}\}_{i=1}^{M}$. Bước cập nhật của thuật
toán hạ gradient là
\begin{equation}
\btheta_{n+1}
= \btheta_{n}
- \eta \left[
    \nabla_{\btheta} J_{r}(\btheta_{n})
    + \sum_{i=1}^{M} \lambda_{i}\, \nabla_{\btheta} J_{i}(\btheta_{n})
  \right].
\label{eq:gd-update}
\end{equation}
 
Hướng cập nhật là \emph{tổng} của các gradient thành phần. Nếu các gradient này
có độ lớn chênh lệch nhau nhiều bậc, thành phần trội chi phối hoàn toàn quỹ đạo
tối ưu, còn các thành phần còn lại bị bỏ qua trên thực tế. Đó là bản chất của
bệnh lý gradient. Cần thấy rằng đây không phải hiện tượng riêng của PINNs mà là
đặc tính chung của mọi phương pháp phạt; điều đặc thù ở đây là \emph{nguyên nhân}
gây mất cân bằng, được truy nguyên ở hai mục sau.
 
\subsection{Dòng gradient và tính cứng của hệ}
\label{subsec:dong-gradient}
 
Xét giới hạn liên tục của \eqref{eq:gd-update} khi tốc độ học $\eta \to 0$, ta
thu được hệ phương trình vi phân thường mô tả \emph{dòng gradient}:
\begin{equation}
\frac{d\btheta}{ds} = -\nabla_{\btheta} J\big(\btheta(s)\big),
\qquad \btheta(0) = \btheta_{0}.
\label{eq:gradient-flow}
\end{equation}
Hạ gradient với bước $\eta$ chính là lược đồ Euler hiện áp dụng cho
\eqref{eq:gradient-flow}, nên các kết luận cổ điển về ổn định của Euler hiện áp
dụng được trực tiếp. Tuyến tính hoá quanh một điểm dừng $\btheta^{\star}$ cho
\begin{equation}
\frac{d}{ds}\,\delta\btheta
\approx -\, \bm{H}\, \delta\btheta,
\qquad
\bm{H} = \nabla^{2}_{\btheta} J(\btheta^{\star}),
\label{eq:linearized-flow}
\end{equation}
với nghiệm phân rã theo các trị riêng $\mu_{1} \ge \cdots \ge \mu_{P} \ge 0$ của
$\bm{H}$. Trên không gian con sinh bởi các hướng riêng có $\mu_i > 0$, điều kiện
ổn định của Euler hiện buộc $\eta < 2/\mu_{1}$, trong khi tốc độ hội tụ theo
hướng riêng chậm nhất bị chi phối bởi trị riêng dương nhỏ nhất $\mu_{\min}^{+}$.
Số vòng lặp cần thiết vì vậy tỉ lệ với số điều kiện hiệu dụng
$\kappa_{\mathrm{eff}} = \mu_{1}/\mu_{\min}^{+}$.
 
\begin{nhanxet}[Vì sao phải viết $\mu_{\min}^{+}$ chứ không phải $\mu_P$]
Với mạng dư thừa tham số, $\bm{H}$ tại một cực tiểu nói chung \emph{suy biến}:
các hướng riêng ứng với trị riêng $0$ là những hướng trong không gian tham số mà
dọc theo đó hàm biểu diễn không đổi. Nhận xét~\ref{nx:doi-xung-hoan-vi} đã cho
một nguồn sinh ra chúng. Dọc các hướng ấy, dòng \eqref{eq:gradient-flow} đứng
yên và không gây khó khăn gì; số điều kiện có ý nghĩa do đó phải được lấy trên
phần bù trực giao của nhân $\bm{H}$. Viết $\kappa = \mu_1/\mu_P$ như một số tài
liệu là không đúng khi $\mu_P = 0$.
\end{nhanxet}
 
Wang và cộng sự chỉ ra rằng với PINNs, hệ \eqref{eq:gradient-flow} là một hệ
\emph{cứng} (stiff): phổ của $\bm{H}$ trải rộng qua nhiều bậc độ lớn, và độ cứng
ấy bị chi phối bởi thành phần phần dư. Bằng chứng số trên bài toán Helmholtz
trong~\cite{wang2021} cho thấy nhiều trị riêng của $\nabla^2_{\btheta}J_r$ đạt
tới cỡ $10^{5}$, trong khi phổ của $\nabla^2_{\btheta}J_b$ vẫn ở thang nhỏ; hệ
quả quan sát được là $J_r$ dao động mạnh ngay cả khi dùng hạ gradient toàn phần,
còn $J_b$ giảm đều. Nguồn gốc giải tích của hiện tượng này nằm ở chính cấu trúc
của $J_r$, vốn chứa đạo hàm bậc cao của mạng theo biến đầu vào; ta quay lại ở
Mục~\ref{subsec:he-qua-pinn}.
 
\subsection{Biểu hiện và hệ quả của bệnh lý gradient}
\label{subsec:bieu-hien}
 
Quan sát thực nghiệm cốt lõi trong~\cite{wang2021}: khi khảo sát phân bố độ lớn
của các gradient lan truyền ngược tại từng lớp ẩn, phân bố của
$\nabla_{\btheta} J_{r}$ tập trung ở những giá trị lớn hơn nhiều bậc so với phân
bố của $\lambda_{i} \nabla_{\btheta} J_{i}$. Hiện tượng trầm trọng hơn khi độ sâu
mạng tăng -- phù hợp với Mệnh đề~\ref{md:tich-gradient}, theo đó tín hiệu ngược
đi qua nhiều lớp bị nhân với một tích các thừa số nhỏ hơn $1$ -- và khi bài toán
có tính cứng cao, chẳng hạn độ nhớt nhỏ trong phương trình Burgers.
 
Hệ quả nghiêm trọng hơn vẻ ngoài kỹ thuật của vấn đề, và có thể phát biểu dưới
góc nhìn tính đặt chỉnh như sau.
 
\begin{nhanxet}[Vì sao gradient biên bị lấn át lại dẫn đến nghiệm sai]
\label{rem:ill-posedness}
Phương trình \eqref{eq:pde-general} khi tách rời khỏi các điều kiện
\eqref{eq:ic}--\eqref{eq:bc} là bài toán \emph{đặt không chỉnh}: tập nghiệm của
nó là một họ vô hạn hàm. Tương ứng, tập
\begin{equation}
\mathcal{M}_{r}
 = \left\{ \btheta \in \R^{P} : J_{r}(\btheta) = 0 \right\}
\label{eq:manifold-r}
\end{equation}
nói chung là một tập có số chiều lớn trong không gian tham số. Vai trò của các
thành phần $J_{i}$ là \emph{chọn ra} một điểm trên tập đó, tương ứng với nghiệm
vật lý.
 
Khi $\nabla_{\btheta} J_{i}$ bị lấn át, quá trình tối ưu vẫn giảm $J_{r}$ rất
hiệu quả và tiến về $\mathcal{M}_{r}$, nhưng trôi tự do trên đó. Mạng thu được
thoả mãn phương trình vi phân với sai số nhỏ song vi phạm điều kiện biên: nó xấp
xỉ \emph{một nghiệm nào đó} của phương trình chứ không phải nghiệm ta cần. Điều
nguy hiểm là giá trị hàm mục tiêu tổng vẫn có thể trông rất nhỏ, khiến người dùng
lầm tưởng mô hình đã hội tụ tốt.
 
Nối với Mục~\ref{subsec:chan-on-dinh}: đây chính là trường hợp giả thiết của
Mệnh đề~\ref{md:chan-on-dinh} bị vi phạm nặng. Phần dư nhỏ không còn kéo theo sai
số nhỏ, vì chặn ổn định đòi hỏi điều kiện biên được thoả. Bởi vậy việc theo dõi
\emph{riêng rẽ} từng thành phần $J_{r}, J_{0}, J_{b}$ trong quá trình huấn luyện
là bắt buộc, và ta áp dụng nguyên tắc này trong toàn bộ thực nghiệm ở
Chương~\ref{chap:thucnghiem}. Một giá trị $J$ tổng nhỏ, tự nó, không phải bằng
chứng của hội tụ.
\end{nhanxet}
 
\subsection{Cân bằng tốc độ học thích nghi}
\label{subsec:annealing}
 
Trước khi trình bày thuật toán, cần trả lời một câu hỏi hiển nhiên: nếu
$\lambda_i$ là siêu tham số khó chọn, sao không để thuật toán tự học chúng bằng
chính hạ gradient?
 
\begin{bode}[Không thể học $\lambda_i$ bằng cực tiểu hoá $J$]
\label{bd:lambda-tam-thuong}
Xét $J$ ở dạng \eqref{eq:loss-compact} như một hàm của cặp
$(\btheta, \bm{\lambda})$ với $\bm{\lambda} \in \R^{M}_{>0}$. Khi đó
\begin{equation}
\frac{\partial J}{\partial \lambda_{i}} = J_{i}(\btheta) \;\ge\; 0 ,
\label{eq:dJ-dlambda}
\end{equation}
nên hạ gradient theo $\lambda_i$ là đơn điệu không tăng, và nếu
$J_i(\btheta) > 0$ thì nó đẩy $\lambda_{i}$ về $0$. Cực tiểu của $J$ theo
$\bm{\lambda}$ trên $\R^M_{\ge 0}$ đạt tại $\bm{\lambda} = \bm{0}$, tức tại đó
mọi ràng buộc bị loại bỏ hoàn toàn.
\end{bode}
 
\begin{proof}
$J$ affine theo từng $\lambda_i$ với hệ số $J_i(\btheta)$, cho
\eqref{eq:dJ-dlambda}; các $J_i$ là trung bình của các bình phương nên không âm.
Bước cập nhật $\lambda_i \gets \lambda_i - \eta J_i$ giảm ngặt khi $J_i > 0$.
\end{proof}
 
Bổ đề~\ref{bd:lambda-tam-thuong} nói rằng trọng số cân bằng \emph{không thể} là
biến tối ưu của cùng bài toán: mọi ràng buộc đều làm hàm mục tiêu tăng, nên cách
rẻ nhất để giảm $J$ là bỏ ràng buộc. Chúng phải được xác định bằng một quy tắc
bên ngoài. Wang và cộng sự~\cite{wang2021} đề xuất một quy tắc dựa trên thống kê
gradient, gọi là \emph{cân bằng tốc độ học} (learning rate annealing): chọn
$\lambda_{i}$ sao cho độ lớn đặc trưng của $\lambda_{i}\nabla_{\btheta} J_{i}$
tương đương với độ lớn đặc trưng của $\nabla_{\btheta} J_{r}$.
 
Cụ thể, tại vòng lặp thứ $n$ ta tính giá trị đề xuất
\begin{equation}
\hat{\lambda}_{i}
= \frac{\norm{\nabla_{\btheta} J_{r}(\btheta_{n})}_{\infty}}
       {\dfrac{1}{P}\norm{\nabla_{\btheta} J_{i}(\btheta_{n})}_{1}} ,
\label{eq:lambda-hat}
\end{equation}
trong đó tử số là giá trị tuyệt đối \emph{lớn nhất trong các thành phần} của
vectơ gradient phần dư, còn mẫu số là trung bình cộng các giá trị tuyệt đối của
các thành phần gradient thứ $i$. Cần đọc \eqref{eq:lambda-hat} đúng: cả hai phép
lấy cực đại và trung bình đều thực hiện \emph{trên các thành phần của một vectơ
trong $\R^P$}, không phải trên biến $\btheta$. Việc dùng chuẩn vô cùng ở tử số và
chuẩn $L^1$ chuẩn hoá ở mẫu số là một lựa chọn bất đối xứng có chủ ý: nó cân bằng
\emph{đỉnh} của gradient phần dư với \emph{mức nền} của gradient ràng buộc, cho
$\hat\lambda_i$ lớn hơn so với việc so hai trung bình.
 
Trọng số thực tế được cập nhật theo trung bình trượt để giảm dao động:
\begin{equation}
\lambda_{i}^{(n+1)}
= (1-\alpha)\, \lambda_{i}^{(n)}
+ \alpha\, \hat{\lambda}_{i},
\qquad \alpha \in (0,1).
\label{eq:lambda-update}
\end{equation}
Bài báo gốc khuyến nghị $\alpha = 0{,}1$ và báo cáo độ nhạy rất thấp với tham số
này trong khoảng $[0{,}05;\,0{,}2]$. Giá trị nhỏ của $\alpha$ là điều kiện để
\eqref{eq:lambda-update} thực sự làm trơn: hằng số thời gian của trung bình trượt
là $1/\alpha \approx 10$ lần cập nhật, trong khi $\alpha$ gần $1$ sẽ khiến
$\lambda_i^{(n+1)} \approx \hat\lambda_i$ và mọi tác dụng làm trơn biến mất.
 
\begin{algorithm}[htbp]
\caption{Huấn luyện PINNs với cân bằng tốc độ học thích nghi}
\label{alg:annealing}
\begin{algorithmic}[1]
\Require Kiến trúc $u_{\btheta}$; tốc độ học $\eta = 10^{-3}$; hệ số trung bình
  trượt $\alpha = 0{,}1$; chu kỳ cập nhật $N_{\text{up}}$; số vòng $N_{\max}$
\State Khởi tạo $\btheta_{0}$ theo Xavier \eqref{eq:xavier}; đặt
  $\lambda_{i}^{(0)} = 1$ với mọi $i$
\For{$n = 0, 1, \dots, N_{\max}-1$}
  \If{$n \bmod N_{\text{up}} = 0$}
    \State Tính riêng rẽ $\nabla_{\btheta}J_{r}(\btheta_{n})$ và
           $\nabla_{\btheta}J_{i}(\btheta_{n})$, $i=1,\dots,M$
    \State $\hat{\lambda}_{i} \gets
           \norm{\nabla_{\btheta}J_{r}(\btheta_{n})}_{\infty}
           \big/ \big(\tfrac{1}{P}\norm{\nabla_{\btheta}J_{i}
           (\btheta_{n})}_{1}\big)$
           \Comment{\eqref{eq:lambda-hat}, tách khỏi đồ thị}
    \State $\lambda_{i}^{(n+1)} \gets
           (1-\alpha)\lambda_{i}^{(n)} + \alpha\hat{\lambda}_{i}$
  \Else
    \State $\lambda_{i}^{(n+1)} \gets \lambda_{i}^{(n)}$
  \EndIf
  \State $\btheta_{n+1} \gets \btheta_{n} - \eta\,\nabla_{\btheta}
         \Big[J_{r}(\btheta_{n})
         + \sum_{i} \lambda_{i}^{(n+1)} J_{i}(\btheta_{n})\Big]$
\EndFor
\State \Return $\btheta_{N_{\max}}$
\end{algorithmic}
\end{algorithm}
 
Hai chi tiết cài đặt cần nói rõ, vì cả hai đều là nguồn lỗi thường gặp.
 
\emph{Thứ nhất, phải tách $\hat\lambda_i$ khỏi đồ thị tính toán.} Đại lượng
$\hat{\lambda}_{i}$ trong \eqref{eq:lambda-hat} được tính \emph{từ} các gradient,
nên bản thân nó là một hàm của $\btheta_n$ và nằm trong đồ thị. Nếu không tách
(trong PyTorch: \texttt{detach}), lượt ngược ở dòng 10 sẽ chảy qua $\hat\lambda_i$
và sinh ra các số hạng đạo hàm bậc hai của $J$ theo $\btheta$ -- vừa không nằm
trong ý định của thuật toán, vừa làm đồ thị dài thêm đáng kể theo tinh thần
Nhận xét~\ref{nx:chi-phi-bac-cao}. Cần phân biệt điều này với
Bổ đề~\ref{bd:lambda-tam-thuong}: bổ đề nói vì sao $\lambda_i$ không được là
\emph{tham số học}, còn ở đây ta nói vì sao nó cũng không được là \emph{đỉnh có
gradient} trong đồ thị. Hai lỗi khác nhau, cùng dẫn tới hành vi sai.
 
\emph{Thứ hai, chu kỳ cập nhật.} Tính riêng rẽ gradient của từng thành phần đắt
hơn tính gradient của tổng: theo Mệnh đề~\ref{md:chi-phi-bp}, mỗi lượt ngược tốn
xấp xỉ chi phí một lượt xuôi, nên $M+1$ lượt riêng rẽ đắt gấp $M+1$ lần. Đây là
lý do chỉ cập nhật sau mỗi $N_{\text{up}}$ vòng lặp (thường $10$ hoặc $100$),
khiến chi phí phụ trội chia đều ra thành không đáng kể.
 
\subsection{Kiến trúc mạng có cổng}
\label{subsec:gated}
 
Song song với cơ chế cân bằng trọng số, Wang và cộng sự đề xuất một biến thể kiến
trúc nhằm giảm nhẹ hiện tượng gradient suy giảm theo độ sâu. Kiến trúc bổ sung
hai đường mã hoá $U$ và $V$ được tính \emph{một lần} từ đầu vào, sau đó trộn vào
mọi lớp ẩn qua một cơ chế cổng:
\begin{align}
U &= \sigma\!\left(X W^{(1)} + b^{(1)}\right),
\qquad
V = \sigma\!\left(X W^{(2)} + b^{(2)}\right),
\label{eq:encoders}\\[3pt]
H^{(1)} &= \sigma\!\left(X W^{h} + b^{h}\right),
\label{eq:first-hidden}\\[3pt]
Z^{(k)} &= \sigma\!\left(H^{(k)} W^{z,k} + b^{z,k}\right),
&& k = 1, \dots, L,
\label{eq:gate}\\[3pt]
H^{(k+1)} &= \left(1 - Z^{(k)}\right) \odot U + Z^{(k)} \odot V,
&& k = 1, \dots, L,
\label{eq:mix}\\[3pt]
u_{\btheta}(X) &= H^{(L+1)} W + b,
\label{eq:output-improved}
\end{align}
trong đó $\odot$ là tích Hadamard và $\sigma = \tanh$. Bộ trọng số $W^{h}, b^{h}$
của lớp ẩn đầu tiên là riêng biệt, không dùng chung với $W^{z,1}, b^{z,1}$ của
cổng thứ nhất; hai phép nhân có kích thước khác nhau.
 
Cơ chế \eqref{eq:mix} tạo đường tắt trực tiếp từ đầu vào tới mọi lớp ẩn, theo
tinh thần kết nối dư của mạng ResNet. Lập luận về vì sao nó giúp: theo
Mệnh đề~\ref{md:tich-gradient}, tín hiệu ngược đi từ lớp $L$ về lớp $1$ bị nhân
với một tích $L-1$ ma trận, và tích ấy suy giảm theo cấp số nhân. Trong
\eqref{eq:mix}, $U$ và $V$ nhận gradient \emph{trực tiếp} từ mọi lớp $k$ chứ
không phải qua chuỗi dài, nên đường ngắn nhất từ hàm mục tiêu về $W^{(1)},
W^{(2)}$ có độ dài $\mathcal{O}(1)$ thay vì $\mathcal{O}(L)$. Hình~\ref{fig:gated-arch}
thể hiện cấu trúc này.
 
\begin{figure}[htbp]
\centering
\resizebox{\textwidth}{!}{%
\begin{tikzpicture}[
  font=\small,
  blk/.style={rounded corners=2pt, draw=black!60, line width=0.6pt,
              align=center, inner sep=4.5pt, minimum height=8mm, fill=white},
  enc/.style={blk, fill=orange!12},
  hid/.style={blk, fill=blue!7},
  gate/.style={blk, fill=green!9},
  mix/.style={circle, draw=black!65, line width=0.7pt, minimum size=7mm,
              inner sep=0pt, fill=yellow!14},
  flow/.style={-{Stealth[length=4.5pt]}, line width=0.7pt, black!70},
  side/.style={-{Stealth[length=4pt]}, line width=0.6pt, orange!65!black},
]
\node[blk, fill=black!5] (X) at (0,0) {$X=(x,t)$};
\node[enc] (U) at (3.1, 2.3) {$U=\sigma\!\left(XW^{(1)}+b^{(1)}\right)$};
\node[enc] (V) at (3.1,-2.3) {$V=\sigma\!\left(XW^{(2)}+b^{(2)}\right)$};
\draw[flow] (X.north) |- (U.west);
\draw[flow] (X.south) |- (V.west);
\node[hid]  (H1) at (3.1, 0) {$H^{(1)}$};
\node[gate] (Z1) at (5.6, 0) {$Z^{(1)}$};
\node[mix]  (M1) at (7.7, 0) {};
\node[hid]  (H2) at (9.5, 0) {$H^{(2)}$};
\node       (dots) at (11.0,0) {$\cdots$};
\node[gate] (ZL) at (12.5, 0) {$Z^{(L)}$};
\node[mix]  (ML) at (14.6, 0) {};
\node[hid]  (HL) at (16.5, 0) {$H^{(L+1)}$};
\node[blk, fill=black!5] (OUT) at (19.1, 0) {$u_{\btheta} = H^{(L+1)}W+b$};
\draw[flow] (X)   -- (H1);   \draw[flow] (H1)  -- (Z1);
\draw[flow] (Z1)  -- (M1);   \draw[flow] (M1)  -- (H2);
\draw[flow] (H2)  -- (dots); \draw[flow] (dots)-- (ZL);
\draw[flow] (ZL)  -- (ML);   \draw[flow] (ML)  -- (HL);
\draw[flow] (HL)  -- (OUT);
\node at (M1) {\scriptsize $\odot$};
\node at (ML) {\scriptsize $\odot$};
\draw[side] (U.south) -- (M1.north);
\draw[side] (V.north) -- (M1.south);
\draw[side] (U.east) -- ++(0.7,0) -| (ML.north);
\draw[side] (V.east) -- ++(0.7,0) -| (ML.south);
\node[align=center, font=\footnotesize, black!75] (rule) at (11.0,-3.15)
  {quy tắc trộn tại mỗi nút $\odot$:\quad
   $H^{(k+1)}=\left(1-Z^{(k)}\right)\odot U + Z^{(k)}\odot V$};
\draw[decorate, decoration={brace, amplitude=4pt, mirror},
      draw=black!45, line width=0.5pt]
  ($(M1.south)+(-0.45,-0.75)$) -- ($(ML.south)+(0.45,-0.75)$);
\node[font=\footnotesize, orange!65!black, align=center] at (3.1,3.35)
  {hai bộ mã hoá dùng chung, tính một lần từ đầu vào};
\end{tikzpicture}}
\caption[Kiến trúc truyền thẳng có cổng]{Kiến trúc truyền thẳng có cổng theo đề
xuất của Wang và cộng sự~\cite{wang2021}. Hai bộ mã hoá $U$ và $V$ được tính một
lần duy nhất từ đầu vào rồi trộn vào mọi lớp ẩn qua cổng $Z^{(k)}$ theo công thức
\eqref{eq:mix}. Các đường nối màu cam là đường tắt nối trực tiếp đầu vào với từng
lớp sâu, rút ngắn quãng đường lan truyền ngược tới $W^{(1)}, W^{(2)}$ từ
$\mathcal{O}(L)$ xuống $\mathcal{O}(1)$.}
\label{fig:gated-arch}
\end{figure}
 
\begin{nhanxet}
\label{nx:hai-huong-bo-tro}
Cơ chế cân bằng trọng số và kiến trúc có cổng giải quyết hai khía cạnh khác nhau
của cùng một vấn đề và có thể kết hợp: cơ chế thứ nhất tác động lên \emph{tỉ lệ}
giữa các gradient thành phần, cơ chế thứ hai tác động lên \emph{đường truyền} của
gradient qua độ sâu. Trên bài toán Helmholtz, \cite{wang2021} báo cáo sai số
tương đối $L^2$ giảm từ $1{,}81\times10^{-1}$ (PINN gốc) xuống
$1{,}27\times10^{-2}$ khi dùng cân bằng thích nghi -- cải thiện hơn một bậc độ
lớn.
 
Cần giữ đúng phạm vi của kết luận: cả hai kỹ thuật đều tác động ở tầng
\emph{tối ưu}. Không kỹ thuật nào trong số đó thay đổi hằng số ổn định của
Mệnh đề~\ref{md:chan-suy-bien}, cũng không thay đổi khả năng biểu diễn tần số cao
của lớp hàm -- nguyên nhân gốc rễ mà ta phân tích ngay sau đây.
\end{nhanxet}
 
%==============================================================================
\section{Thiên kiến phổ}
\label{sec:spectral-bias}
%==============================================================================
 
Mục~\ref{sec:gradient-pathology} phân tích khó khăn của PINNs từ góc độ tối ưu.
Mục này đi sâu hơn một tầng, tới góc độ \emph{lý thuyết xấp xỉ}: ngay cả khi
thuật toán tối ưu hoạt động hoàn hảo, bản thân mạng nơ-ron truyền thẳng vẫn có
một thiên hướng nội tại khiến nó học các thành phần tần số cao rất chậm. Hiện
tượng này được Rahaman và cộng sự~\cite{rahaman2019} gọi là \emph{thiên kiến phổ}
(spectral bias).
 
\subsection{Đặt vấn đề: thang độ dài nhỏ và tần số cao}
\label{subsec:thang-do-dai}
 
Nhiều bài toán vi phân có nghiệm chứa cấu trúc biến thiên nhanh trên một thang độ
dài rất nhỏ. Hai lớp điển hình, cả hai đều đã xuất hiện ở
Mục~\ref{subsec:chan-on-dinh}.
 
\emph{Lớp biên trong bài toán nhiễu kỳ dị.} Phương trình $-\nu u'' + u' = f$ trên
$[0,1]$ với $0 < \nu \ll 1$ -- chính là bài toán của
Mệnh đề~\ref{md:chan-suy-bien} -- có nghiệm chứa một lớp biên bề dày
$\mathcal{O}(\nu)$ gần $x = 1$.
 
\emph{Lớp trong của phương trình Burgers.} Với điều kiện đầu
$u(x,0) = -\sin(\pi x)$ trên $[-1,1]$ và độ nhớt
$\nu = 0{,}01/\pi \approx 3{,}18 \times 10^{-3}$, nghiệm phát triển một lớp
chuyển tiếp dốc quanh $x = 0$. Nghiệm tham chiếu tính bằng phương pháp giả phổ
Fourier cho thấy độ dốc đạt cực đại $\max\abs{u_x} \approx 151$ tại
$t \approx 0{,}51$, tương ứng bề dày lớp
\begin{equation}
\delta := \frac{2 \max_{x} \abs{u}}{\max_{x} \abs{u_x}}
\approx 1{,}3 \times 10^{-2} \approx 4\nu .
\label{eq:layer-width}
\end{equation}
 
Theo nguyên lý bất định trong giải tích Fourier, một cấu trúc có bề rộng $\delta$
trong miền không gian tương ứng với các thành phần phổ có tần số cỡ
$\abs{k} \sim 1/\delta$. Với lớp trong của Burgers, từ \eqref{eq:layer-width} ta
được $\abs{k} \sim 77$: các tần số cỡ $10^{2}$ đóng vai trò quyết định, cao hơn
hai bậc so với tần số $\abs{k} = 1$ của điều kiện đầu $-\sin(\pi x)$. Câu hỏi đặt
ra là: một mạng nơ-ron truyền thẳng biểu diễn các tần số như vậy tốt đến đâu?
 
\begin{figure}[htbp]
\centering
\includegraphics[width=\textwidth]{Images/chap_4/hinh_burgers.pdf}
\caption[Lớp trong của nghiệm phương trình Burgers]{Nghiệm tham chiếu của phương
trình Burgers với $\nu = 0{,}01/\pi$, tính bằng phương pháp giả phổ Fourier
($N = 2048$ chế độ, tích phân thời gian RK4 với thừa số tích phân). Bảng (a): bản
đồ nhiệt của $u(x,t)$; lớp chuyển tiếp hình thành quanh $x=0$ và đạt độ dốc lớn
nhất khoảng $t \approx 0{,}51$. Bảng (b): lát cắt tại ba thời điểm, kèm khung
phóng to quanh $x = 0$. Ở tỉ lệ đầy đủ, lớp trong hầu như không nhìn thấy được --
chính đặc điểm này khiến nó đòi hỏi các thành phần phổ tần số cao, vốn là những
thành phần mà mạng nơ-ron học chậm nhất (Định lý~\ref{dl:spectral-decay}).}
\label{fig:burgers}
\end{figure}
 
\subsection{Phân tích Fourier của mạng ReLU}
\label{subsec:fourier-relu}
 
Rahaman và cộng sự tiếp cận câu hỏi trên bằng cách phân tích trực tiếp phổ Fourier
của hàm mà mạng biểu diễn. Điểm khởi đầu là một tính chất cấu trúc đã được chứng
minh ở Chương~\ref{chap:fnn}.
 
\begin{bode}[Biểu diễn tuyến tính từng khúc]
\label{bd:cpwl}
Mọi mạng truyền thẳng $f_{\btheta} : \R^{d} \to \R$ dùng hàm kích hoạt
$\mathrm{ReLU}$ đều là hàm tuyến tính liên tục từng khúc: tồn tại phân hoạch hữu
hạn miền xác định thành các đa diện lồi $\{P_{i}\}_{i=1}^{R}$ sao cho
\begin{equation}
f_{\btheta}(x) = \sum_{i=1}^{R}
\mathbf{1}_{P_{i}}(x) \left( \bm{W}_{i}\, x + b_{i} \right),
\label{eq:cpwl}
\end{equation}
trong đó $\mathbf{1}_{P_{i}}$ là hàm chỉ báo của $P_{i}$, còn $\bm{W}_{i}, b_{i}$
là các tích tương ứng của các ma trận trọng số theo mẫu kích hoạt trên $P_{i}$.
\end{bode}
 
Đây chính là Định lý~\ref{dl:relu-tuyen-tinh}, phát biểu lại dưới dạng công thức
tường minh; chứng minh đã cho ở Mục~\ref{subsec:relu}.
 
Từ biểu diễn \eqref{eq:cpwl}, phổ Fourier được ước lượng bằng cách khai thác
quan hệ giữa đạo hàm và phép nhân trong miền tần số. Nếu $\widetilde{f}$ là biến
đổi Fourier của $f$ thì $\widetilde{\partial^{\alpha} f}(k)
= (2\pi i k)^{\alpha}\widetilde{f}(k)$, nên
$\widetilde{f}(k) = \widetilde{\partial^{\alpha}f}(k)/(2\pi i k)^{\alpha}$: mỗi
bậc đạo hàm mà ta ``lấy ra được'' của $f$ cho thêm một thừa số $1/\abs{k}$ trong
ước lượng phổ. Với $f_{\btheta}$ tuyến tính từng khúc, sau $d+1$ bậc thì đạo hàm
suy rộng trở thành một tổ hợp các phân bố delta trên các mặt của phân hoạch --
đúng hiện tượng đã gặp ở Mệnh đề~\ref{md:dirac} cho trường hợp một chiều -- và
biến đổi Fourier của một tổ hợp hữu hạn các delta thì bị chặn đều. Kết quả dẫn
tới định lý sau.
 
\begin{dinhly}[Suy giảm phổ của mạng ReLU~\cite{rahaman2019}]
\label{dl:spectral-decay}
Cho $f_{\btheta}$ là mạng ReLU như ở Bổ đề~\ref{bd:cpwl}, xét trên một miền bị
chặn. Biến đổi Fourier $\widetilde{f}_{\btheta}(k)$ của nó thoả mãn
\begin{equation}
\abs{\widetilde{f}_{\btheta}(k)}
\;\le\;
\frac{C\, L_{\btheta}}{\norm{k}^{\,d+1}},
\label{eq:fourier-decay}
\end{equation}
trong đó $C$ chỉ phụ thuộc miền xác định và số mảnh của phân hoạch, còn
$L_{\btheta}$ bị chặn bởi một đa thức theo chuẩn của các ma trận trọng số
$\{\bW^{(\ell)}\}$.
\end{dinhly}
 
\begin{nhanxet}[Ba điều \eqref{eq:fourier-decay} nói và một điều nó không nói]
\label{nx:doc-dinh-ly-pho}
Việc đọc đúng định lý này quan trọng, vì nó thường được diễn giải quá mạnh.
 
Thứ nhất, giả thiết miền bị chặn là cần thiết: một hàm tuyến tính từng khúc trên
toàn $\R^d$ không khả tích, nên biến đổi Fourier chỉ có nghĩa theo nghĩa phân bố
hoặc sau khi hạn chế miền.
 
Thứ hai, \eqref{eq:fourier-decay} là một \emph{chặn trên}. Tự nó, nó không chứng
minh rằng mạng \emph{không thể} biểu diễn tần số cao; nó chỉ ràng buộc biên độ
phổ với chuẩn tham số. Sức nặng của lập luận nằm ở việc đảo ngược quan hệ ấy, tức
Hệ quả~\ref{hq:parameter-cost}.
 
Thứ ba, định lý phát biểu cho ReLU. Mạng $\tanh$ dùng trong báo cáo không thuộc
phạm vi của nó, và thực tế phổ của mạng $\tanh$ suy giảm còn nhanh hơn -- theo
Mệnh đề~\ref{md:tanh-giai-tich}, $\tanh$ giải tích trên một dải, nên hàm hợp
tương ứng trơn hơn hẳn tuyến tính từng khúc. Kết luận định tính (tần số cao đắt
đỏ) chuyển sang được; hằng số và số mũ cụ thể thì không. Ta dẫn
Định lý~\ref{dl:spectral-decay} như bằng chứng cho cơ chế, không như một chặn áp
dụng trực tiếp cho kiến trúc của báo cáo.
\end{nhanxet}
 
\begin{hequa}[Cái giá tham số của tần số cao]
\label{hq:parameter-cost}
Để $f_{\btheta}$ biểu diễn được một thành phần phổ có biên độ $A$ tại tần số $k$,
chuẩn tham số phải thoả mãn
\begin{equation}
L_{\btheta} \;\ge\; \frac{A}{C}\,\norm{k}^{\,d+1}.
\label{eq:param-cost}
\end{equation}
\end{hequa}
 
\begin{proof}
Nếu $\abs{\widetilde f_{\btheta}(k)} = A$ thì \eqref{eq:fourier-decay} cho
$A \le C L_{\btheta}/\norm{k}^{d+1}$; chuyển vế.
\end{proof}
 
Tần số càng cao thì mạng càng buộc phải dùng trọng số lớn, và cái giá tăng theo
luỹ thừa $d+1$. Điều này có hai hệ luỵ thực hành, cả hai nối trực tiếp với
Chương~\ref{chap:fnn}.
 
Thứ nhất, khởi tạo Xavier \eqref{eq:xavier} cho chuẩn trọng số nhỏ theo
thiết kế -- đó chính là mục đích của nó, giữ $\bz^{(\ell)}$ trong vùng gần tuyến
tính của $\tanh$ để tránh bão hoà (Mệnh đề~\ref{md:tich-gradient}). Nhưng
\eqref{eq:param-cost} nói rằng biểu diễn tần số cao đòi hỏi chuẩn trọng số
\emph{lớn}. Hai yêu cầu này xung đột: quá trình huấn luyện phải đi một quãng
đường dài trong không gian tham số, xuất phát từ vùng ``an toàn'' về gradient
hướng tới vùng ``đắt'' về chuẩn.
 
Thứ hai, mọi cơ chế chính quy hoá dựa trên chuẩn tham số -- suy giảm trọng số,
dừng sớm -- đều \emph{trực tiếp cản trở} việc học tần số cao. Với bài toán có lớp
biên, chính quy hoá không chỉ vô ích mà còn phản tác dụng.
 
\subsection{Thiên kiến trong động lực học huấn luyện}
\label{subsec:dong-luc-hoc}
 
Định lý~\ref{dl:spectral-decay} nói về khả năng biểu diễn tại một bộ tham số cố
định. Rahaman và cộng sự bổ sung một kết quả thứ hai về \emph{quá trình} học: các
thành phần tần số thấp không chỉ dễ biểu diễn hơn mà còn được học \emph{trước},
và bền vững hơn dưới nhiễu loạn tham số. Thực nghiệm của họ trên bài toán hồi quy
hàm nhiều tần số cho thấy các thành phần phổ được khớp gần như tuần tự từ thấp
lên cao.
 
\begin{nhanxet}[Góc nhìn bổ sung qua nhân tiếp tuyến nơ-ron]
\label{rem:ntk-view}
Kết luận trên có thể được củng cố bằng một lập luận độc lập qua nhân tiếp tuyến
nơ-ron (Neural Tangent Kernel). Xét bài toán hồi quy với mất mát bình phương và
dòng gradient \eqref{eq:gradient-flow}. Trong chế độ bề rộng lớn, nơi nhân tiếp
tuyến xấp xỉ không đổi trong suốt quá trình huấn luyện, khai triển sai số
$e(s) = u_{\btheta(s)} - u$ theo hệ hàm riêng $\{\phi_{j}\}$ của toán tử nhân,
với các trị riêng $\varkappa_{1} \ge \varkappa_{2} \ge \cdots \ge 0$, cho
\begin{equation}
\langle e(s), \phi_{j} \rangle
= \langle e(0), \phi_{j} \rangle\, e^{-\varkappa_{j} s}.
\label{eq:ntk-decay}
\end{equation}
Vì phổ $\{\varkappa_{j}\}$ suy giảm nhanh khi tần số của $\phi_{j}$ tăng, các
thành phần sai số tần số cao phân rã với hằng số thời gian lớn hơn nhiều bậc. Đây
chính là thiên kiến phổ nhìn từ động lực học tối ưu, và nó cũng giải thích tính
cứng của hệ \eqref{eq:gradient-flow} đã nêu ở Mục~\ref{subsec:dong-gradient}: tỉ
số $\varkappa_1/\varkappa_{\min}^{+}$ lớn chính là hệ quả của phổ suy giảm nhanh.
 
Hai giả thiết phải nêu rõ. Kết quả \eqref{eq:ntk-decay} đòi hỏi chế độ tuyến tính
hoá (bề rộng lớn, tham số dịch chuyển ít), điều không được bảo đảm cho các mạng
kích thước vừa dùng trong thực nghiệm; và nó là kết quả cho \emph{hồi quy}, chưa
phải cho hàm mục tiêu phần dư. Ta dẫn nó như một lời giải thích cơ chế thống nhất
giữa hai mục, không như một định lý về mô hình đang xét.
\end{nhanxet}
 
\subsection{Hệ quả đối với PINNs: ba nguyên nhân cộng hưởng}
\label{subsec:he-qua-pinn}
 
Với PINNs, thiên kiến phổ nghiêm trọng hơn so với hồi quy thông thường, vì hàm
mục tiêu chứa \emph{đạo hàm} của mạng.
 
\begin{menhde}[Khuếch đại phổ của toán tử vi phân]
\label{md:khuech-dai}
Giả sử $\Lop$ tuyến tính với hệ số hằng và bậc cao nhất $m$, ký hiệu
$\Lop = \sum_{\abs{\alpha}\le m} c_\alpha \partial^{\alpha}$. Với
$e_{\btheta} = u_{\btheta} - u$ và $r_{\btheta} = \Lop[e_{\btheta}]$ (bỏ qua số
hạng thời gian),
\begin{equation}
\widetilde{r}_{\btheta}(k)
 = p(2\pi i k)\, \widetilde{e}_{\btheta}(k),
\qquad
p(\xi) := \sum_{\abs{\alpha}\le m} c_\alpha\, \xi^{\alpha},
\label{eq:residual-amplification}
\end{equation}
nên với các $k$ lớn,
\begin{equation}
\abs{\widetilde{r}_{\btheta}(k)}
 \;\sim\; \abs{p_m\!\left(2\pi i\,k/\norm{k}\right)}\;
   \norm{k}^{m}\,\abs{\widetilde{e}_{\btheta}(k)},
\label{eq:bieu-trung-chinh}
\end{equation}
trong đó $p_m(\xi) := \sum_{\abs{\alpha}=m} c_\alpha \xi^{\alpha}$ là
\emph{biểu trưng chính}. Phát biểu tiệm cận này đòi hỏi
$p_m\big(2\pi i\,k/\norm{k}\big) \ne 0$ theo hướng đang xét --- điều kiện được
thoả với toán tử elliptic, nhưng \emph{không} với mọi toán tử: chẳng hạn toán tử
sóng $\partial_{tt} - \partial_{xx}$ có biểu trưng chính triệt tiêu trên nón ánh
sáng, và ở đúng những hướng ấy không có sự khuếch đại nào.
\end{menhde}
 
\begin{proof}
Biến đổi Fourier biến $\partial^{\alpha}$ thành phép nhân với
$(2\pi i k)^{\alpha}$; lấy tổ hợp tuyến tính với các hệ số $c_\alpha$ cho
\eqref{eq:residual-amplification}. Viết $k = \norm{k}\,\omega$ với
$\norm{\omega} = 1$; do $\xi \mapsto \xi^{\alpha}$ thuần nhất bậc
$\abs{\alpha}$,
\[
p(2\pi i k) = \norm{k}^{m} p_m(2\pi i \omega)
 + \mathcal{O}\big(\norm{k}^{m-1}\big).
\]
Khi $p_m(2\pi i\omega) \ne 0$, số hạng thứ nhất trội và cho
\eqref{eq:bieu-trung-chinh}. Cần nhấn mạnh rằng thừa số
$\abs{p_m(2\pi i\omega)}$ phụ thuộc \emph{hướng} $\omega$ chứ không chỉ phụ
thuộc $\norm{k}$; nó thu về một hằng số duy nhất trong trường hợp một chiều, nơi
chỉ có hai hướng $\omega = \pm 1$ --- và đó là toàn bộ các bài toán được thực
nghiệm ở Chương~\ref{chap:thucnghiem}.
\end{proof}
 
Với $\Lop$ phi tuyến như Burgers, \eqref{eq:residual-amplification} chỉ đúng cho
bài toán đã tuyến tính hoá quanh nghiệm; số hạng $u\partial_x u$ còn sinh thêm
hiện tượng trộn tần số mà phân tích tuyến tính không mô tả được. Ta dùng
Mệnh đề~\ref{md:khuech-dai} như phân tích biểu trưng chính, không như một đẳng thức
cho bài toán đầy đủ.
 
Quan hệ \eqref{eq:residual-amplification} tạo ra một mâu thuẫn mang tính cấu
trúc. Toán tử vi phân \emph{khuếch đại} sai số tần số cao với hệ số
$\norm{k}^{m}$, làm số điều kiện của bài toán tối ưu xấu đi đáng kể -- đây chính
là nguồn gốc giải tích của tính cứng đã mô tả ở \eqref{eq:linearized-flow}, và
giải thích vì sao~\cite{wang2021} quan sát thấy phổ của
$\nabla^2_{\btheta}J_r$ trải tới $10^5$ trong khi phổ của
$\nabla^2_{\btheta}J_b$ vẫn nhỏ: $J_b$ không chứa đạo hàm nào. Mặt khác, theo
Định lý~\ref{dl:spectral-decay} và Hệ quả~\ref{hq:parameter-cost}, mạng lại có xu
hướng nội tại là \emph{không} biểu diễn đúng những tần số ấy.
 
Ghép với Mục~\ref{subsec:chan-on-dinh}, ta thu được bức tranh đầy đủ gồm ba
nguyên nhân độc lập cùng chỉ về một hướng, tóm tắt ở Bảng~\ref{tab:ba-nguyen-nhan}.
 
\begin{table}[htbp]
\centering
\caption{Ba nguyên nhân độc lập khiến PINNs suy giảm độ chính xác trên bài toán
có thang độ dài nhỏ, và tầng phân tích tương ứng.}
\label{tab:ba-nguyen-nhan}
\small
\begin{tabular}{@{}l l p{6.4cm}@{}}
\toprule
\textbf{Tầng} & \textbf{Kết quả} & \textbf{Nội dung} \\
\midrule
Bài toán vi phân & Mệnh đề~\ref{md:chan-suy-bien}
  & Hằng số ổn định tỉ lệ $1/\nu$: cùng một phần dư cho sai số lớn gấp
    $\mathcal{O}(1/\nu)$ lần. \\
\addlinespace
Tối ưu & \eqref{eq:linearized-flow}, \eqref{eq:residual-amplification}
  & Toán tử vi phân khuếch đại tần số cao theo $\norm{k}^m$, làm phổ Hessian
    trải rộng và hệ dòng gradient trở nên cứng. \\
\addlinespace
Xấp xỉ & Hệ quả~\ref{hq:parameter-cost}
  & Biểu diễn tần số $k$ đòi hỏi chuẩn tham số cỡ $\norm{k}^{d+1}$, ngược chiều
    với khởi tạo và với mọi chính quy hoá. \\
\bottomrule
\end{tabular}
\end{table}
 
Ba nguyên nhân này cộng hưởng theo chiều bất lợi: bài toán đòi hỏi độ chính xác
cao nhất ở đúng những tần số mà mô hình học chậm nhất, và ở đúng chế độ tham số
mà chặn ổn định suy biến. Điều này lý giải một quan sát phổ biến trong tài liệu:
PINNs cho kết quả rất tốt với các bài toán trơn có nghiệm biến thiên chậm, nhưng
suy giảm nhanh khi độ nhớt giảm, khi hệ số nhiễu kỳ dị nhỏ đi, hoặc khi nghiệm có
dao động tần số cao; Krishnapriyan và cộng sự~\cite{krishnapriyan2021} khảo sát
có hệ thống đúng hiện tượng này và cho thấy sự suy giảm đến từ \emph{không gian
hàm mục tiêu} chứ không từ khả năng biểu diễn của mạng -- nhất quán với việc
Bảng~\ref{tab:ba-nguyen-nhan} đặt hai trong ba nguyên nhân ở tầng bài toán và
tầng tối ưu. Ta sẽ kiểm chứng định lượng nhận định này ở
Chương~\ref{chap:thucnghiem}, qua cặp đối chứng gồm phương trình khuếch tán
(hệ số $\nu = 0{,}1$, bài toán ``lành mạnh'') và phương trình Burgers
($\nu = 0{,}01/\pi$, bài toán cứng).
 
Cũng cần rút ra một kết luận phương pháp luận từ Bảng~\ref{tab:ba-nguyen-nhan}:
vì ba nguyên nhân nằm ở ba tầng khác nhau, không kỹ thuật nào tác động lên một
tầng có thể sửa được hai tầng còn lại. Cân bằng trọng số
(Thuật toán~\ref{alg:annealing}) chỉ chạm tới tầng giữa; nó không thể cứu một bài
toán mà hằng số ổn định đã suy biến.
 
\subsection{Các hướng khắc phục}
\label{subsec:khac-phuc}
 
Nhận diện được nguyên nhân, ta có thể phân loại các kỹ thuật cải thiện đã được đề
xuất theo cơ chế tác động, như tóm tắt ở Bảng~\ref{tab:remedies}. Điểm chung của
chúng: không kỹ thuật nào loại bỏ thiên kiến phổ, chúng chỉ dịch chuyển hoặc mở
rộng dải tần mà mô hình xử lý hiệu quả.
 
\begin{table}[htbp]
\centering
\caption{Phân loại các kỹ thuật khắc phục thiên kiến phổ trong PINNs.}
\label{tab:remedies}
\small
\begin{tabular}{@{}p{4.0cm} p{4.0cm} p{5.4cm}@{}}
\toprule
\textbf{Kỹ thuật} & \textbf{Cơ chế tác động} & \textbf{Ghi chú} \\
\midrule
Nhúng đặc trưng Fourier &
Đưa sẵn các hàm cơ sở $\sin(2\pi Bx)$, $\cos(2\pi Bx)$ vào đầu vào &
Dịch phổ hiệu dụng của mạng về vùng tần số cao; hiệu quả phụ thuộc mạnh vào việc
chọn thang $B$, vốn đòi hỏi biết trước thang độ dài của nghiệm \\
\addlinespace
Hàm kích hoạt thích nghi &
Thêm hệ số co giãn $a$ khả huấn luyện: $\sigma(a z)$ &
Cho phép mạng tự điều chỉnh thang tần số; chi phí bổ sung không đáng kể; tương
đương một cách tham số hoá lại chuẩn trọng số trong \eqref{eq:param-cost} \\
\addlinespace
Phân rã miền (cPINN, XPINN) &
Chia miền thành các miền con, mỗi miền một mạng riêng &
Mỗi mạng chỉ cần biểu diễn dải tần hẹp; phát sinh điều kiện ghép nối, tức thêm
thành phần mất mát và thêm bệnh lý gradient \\
\addlinespace
Hành quân theo thời gian &
Giải tuần tự trên các khoảng thời gian ngắn &
Tránh phải học lớp trong ngay từ đầu; phù hợp bài toán tiến hoá; sai số tích luỹ
qua các khoảng \\
\addlinespace
Lấy mẫu thích nghi (RAR) &
Bổ sung điểm phối trí tại nơi phần dư lớn &
Tập trung tài nguyên vào vùng có cấu trúc dốc~\cite{lu2021}; giảm sai số lấy mẫu
trong \eqref{eq:mc-estimate} nhưng không đổi lớp hàm \\
\bottomrule
\end{tabular}
\end{table}
 
Đọc Bảng~\ref{tab:remedies} cùng Bảng~\ref{tab:ba-nguyen-nhan} cho một tiêu chí
lựa chọn: nhúng đặc trưng Fourier và hàm kích hoạt thích nghi tác động lên tầng
xấp xỉ; phân rã miền và hành quân theo thời gian tác động lên tầng bài toán bằng
cách thay bài toán khó bằng một dãy bài toán dễ hơn; lấy mẫu thích nghi tác động
lên sai số lấy mẫu, tức một nguồn thứ tư không nằm trong ba nguyên nhân trên. Việc
lựa chọn vì vậy phụ thuộc vào chỗ nào là nút thắt của bài toán cụ thể -- và việc
xác định nút thắt ấy đòi hỏi theo dõi riêng rẽ từng thành phần mất mát cùng sai số
trên tập điểm độc lập, đúng như Nhận xét~\ref{rem:monte-carlo} và
Nhận xét~\ref{rem:ill-posedness} đã lập luận.
 
%==============================================================================
\section*{Kết luận chương}
\addcontentsline{toc}{section}{Kết luận chương}
%==============================================================================
 
Chương này đã trình bày mô hình PINNs theo ba tầng phân tích nối tiếp; các kết
quả chính được tổng hợp ở Bảng~\ref{tab:tong-ket-4}.
 
Ở tầng thứ nhất, ta xây dựng hàm mất mát \eqref{eq:total-loss} từ định nghĩa phần
dư \eqref{eq:residual} và các tập điểm phối trí, đồng thời chỉ ra rằng $J_{r}$ về
bản chất là ước lượng Monte Carlo của chuẩn $L^{2}$ của phần dư -- kèm cảnh báo
rằng ước lượng ấy mất tính không chệch ngay khi $\btheta$ được tối ưu trên chính
tập điểm đó. Quan trọng hơn, Mệnh đề~\ref{md:chan-on-dinh} và
Mệnh đề~\ref{md:chan-suy-bien} thiết lập mắt xích logic mà phương pháp cần: phần
dư nhỏ kéo theo sai số nhỏ khi và chỉ khi bài toán có chặn ổn định, với hằng số
tỉ lệ $1/\nu$. Cùng một khung hàm mất mát phục vụ được cả bài toán thuận lẫn bài
toán ngược, trong đó tham số vật lý được xử lý hoàn toàn như tham số mạng.
 
Ở tầng thứ hai, ta phân tích khó khăn tối ưu phát sinh từ việc gộp nhiều mục tiêu
cạnh tranh. Bệnh lý gradient khiến các thành phần điều kiện biên bị lấn át, và do
bài toán vi phân tách rời khỏi điều kiện biên là đặt không chỉnh, hậu quả là mạng
hội tụ về một nghiệm sai với giá trị hàm mục tiêu trông vẫn nhỏ.
Bổ đề~\ref{bd:lambda-tam-thuong} giải thích vì sao trọng số cân bằng không thể là
biến tối ưu của chính bài toán, do đó phải được xác định bằng một quy tắc bên
ngoài như Thuật toán~\ref{alg:annealing}.
 
Ở tầng thứ ba, ta truy nguyên khó khăn về giới hạn biểu diễn của chính mạng.
Định lý~\ref{dl:spectral-decay} cho thấy phổ Fourier của mạng ReLU suy giảm theo
$\norm{k}^{-(d+1)}$, nên biểu diễn tần số cao đòi hỏi chuẩn tham số lớn -- ngược
chiều với khởi tạo Xavier và với mọi chính quy hoá.
 
Đóng góp tổng hợp của chương là Bảng~\ref{tab:ba-nguyen-nhan}: ba nguyên nhân độc
lập, ở ba tầng khác nhau, cùng làm PINNs suy giảm trên đúng lớp bài toán có thang
độ dài nhỏ. Vì chúng nằm ở ba tầng, không kỹ thuật khắc phục nào tác động lên một
tầng có thể thay thế cho hai tầng còn lại -- một kết luận có hệ quả trực tiếp lên
cách thiết kế và đọc thực nghiệm.
 
\begin{table}[htbp]
\centering
\caption{Tổng kết các kết quả chính của Chương~\ref{chap:pinn}.}
\label{tab:tong-ket-4}
\small
\begin{tabular}{@{}l p{8.2cm}@{}}
\toprule
\textbf{Kết quả} & \textbf{Nội dung và vai trò} \\
\midrule
Định nghĩa~\ref{def:residual}
  & Phần dư dùng chung tham số với mạng; đòi hỏi $u_{\btheta}\in C^m$, nối với
    Hệ quả~\ref{hq:pinn-hong}. \\
Nhận xét~\ref{rem:monte-carlo}
  & $J_r$ là ước lượng Monte Carlo của $\norm{r_{\btheta}}_{L^2}^2$; mất tính
    không chệch khi tối ưu trên chính tập mẫu. \\
\textbf{Mệnh đề~\ref{md:chan-on-dinh}}
  & \textbf{Chặn ổn định} $\norm{e_{\btheta}}_{L^2}\le \pi^{-2}
    \norm{r_{\btheta}}_{L^2}$ cho bài toán Poisson: cơ sở biện minh cho việc cực
    tiểu hoá phần dư. \\
\textbf{Mệnh đề~\ref{md:chan-suy-bien}}
  & \textbf{Hằng số ổn định suy biến như $1/\nu$} với bài toán đối lưu--khuếch
    tán -- tầng thứ nhất của Bảng~\ref{tab:ba-nguyen-nhan}. \\
Bổ đề~\ref{bd:rang-buoc-cung}
  & Nghiệm thử $G + D\,u_{\btheta}$ thoả điều kiện biên với mọi $\btheta$; loại
    $J_b$ khỏi hàm mục tiêu. \\
Nhận xét~\ref{rem:ill-posedness}
  & Gradient biên bị lấn át $\Rightarrow$ trôi trên $\mathcal{M}_r$ $\Rightarrow$
    nghiệm sai với $J$ nhỏ; bắt buộc theo dõi riêng từng thành phần. \\
Bổ đề~\ref{bd:lambda-tam-thuong}
  & $\partial J/\partial\lambda_i = J_i \ge 0$: trọng số cân bằng không học được
    bằng hạ gradient. \\
Định lý~\ref{dl:spectral-decay}, Hệ quả~\ref{hq:parameter-cost}
  & Phổ mạng ReLU suy giảm theo $\norm{k}^{-(d+1)}$; biểu diễn tần số $k$ đòi hỏi
    chuẩn tham số cỡ $\norm{k}^{d+1}$. \\
Mệnh đề~\ref{md:khuech-dai}
  & Toán tử vi phân khuếch đại phổ sai số theo $\norm{k}^{m}$ -- nguồn gốc giải
    tích của tính cứng. \\
\textbf{Bảng~\ref{tab:ba-nguyen-nhan}}
  & \textbf{Ba nguyên nhân độc lập ở ba tầng}, cùng chỉ về một hướng bất lợi. \\
\bottomrule
\end{tabular}
\end{table}
 
Ba tầng phân tích này định hình trực tiếp thiết kế thực nghiệm ở
Chương~\ref{chap:thucnghiem}, nơi ta cài đặt PINNs bằng PyTorch và đánh giá định
lượng cả độ chính xác lẫn hành vi hội tụ của từng thành phần mất mát, trên nhiều
hạt giống ngẫu nhiên theo yêu cầu của Nhận xét~\ref{nx:doi-xung-hoan-vi}.
 